\uuid{EZVx}
\titre{Somme/Produit/DL usuels}
\auteur{hanani ; mzari}
\organisation{exo7}
\licence{CC BY-NC-SA 4.0}
\datecreate{2020-01-20}
\niveau{L1}
\module{Analyse}
\chapitre{Développement limité}
\sousChapitre{Calculs}
\difficulte{1}
\theme{Opérations sur les DL, Somme, Produit, DL usuels}
\exo7qcmid{lille-2.developpements-limites.004}

\contenu{
\texte{
Soit $a\in \Rr$ et $f$ une fonction. On dit que $f(x)$ est un $o\left((x-a)^n\right)$ au voisinage de $a$ et on écrit $f(x)=o((x-a)^n)$ si $\displaystyle \lim _{x\to a}\frac{f(x)}{(x-a)^n}=0$. On note par $DL_n(a)f(x)$ le développement limité de $f$ à l'ordre $n$ en $a$. On note aussi $DL_n(a^+)f(x)$ (resp. $DL_n(a^-)f(x)$) le développement limité de $f$ à l'ordre $n$ à droite en $a$ (resp. à gauche en $a$).
}

\begin{question}
\qclassification{125.02}
\qtags{motcle=Somme/Produit/DL usuels}
Parmi les affirmations suivantes, cocher celles qui sont vraies :
\begin{answers}
\good{$\displaystyle DL_3(0)\sin x=x-\frac{x^3}{6}+o(x^3)$.}
\bad{$\displaystyle DL_3(0)\mathrm{e}^{x}=x+\frac{x^2}{2}+\frac{x^3}{6}+o(x^3)$.}
\bad{$\displaystyle DL_3(0)\left(\sin x+\mathrm{e}^{x}\right)=2x+\frac{x^2}{2}+o(x^3)$.}
\good{$\displaystyle DL_3(0)\left(\sin x\mathrm{e}^{x}\right)=x+x^2+\frac{x^3}{3}+o(x^3)$.}
\end{answers}
\begin{explanations}
On a : $\displaystyle \sin x=x-\frac{x^3}{6}+o(x^3)$ et $\displaystyle \mathrm{e}^{x}=1+x+\frac{x^2}{2}+\frac{x^3}{6}+o(x^3)$. Donc
$$DL_3(0)\left(\mathrm{e}^{x}+\sin x\right)=1+2x+\frac{x^2}{2}+o(x^3)\quad\mbox{et}\quad DL_3(0)\left(\sin x\mathrm{e}^{x}\right)=x+x^2+\frac{x^3}{3}+o(x^3).$$
\end{explanations}
\end{question}
}
